% Auto-generated by scripts/project_calculations.py — do not hand-edit
\begin{table}[htbp]
\centering
\caption{LCOW and annual energy (illustrative SEC and CAPEX from YAML).}
\small
\begin{tabular}{@{}lcccc@{}}
\toprule
Route & SEC [kWh/kg] & $E_{\mathrm{year}}$ [GWh] & CAPEX [MUSD] & LCOW [\$/m$^3$] \\
\midrule
VCR & 0.42 & 15.8233 & 1.80 & 13.98 \\
Absorption & 0.38 & 14.3163 & 1.80 & 13.18 \\
Silica\_adsorption & 1.16 & 43.7023 & 2.20 & 30.02 \\
MOF & 2.96 & 111.5163 & 5.00 & 74.71 \\
TEC\_low & 2 & 75.3488 & 1.60 & 44.96 \\
TEC\_mid & 6 & 226.0465 & 1.60 & 124.96 \\
TEC\_high & 12 & 452.0930 & 1.60 & 244.96 \\
\bottomrule
\end{tabular}
\end{table}

\subsection*{Design Translation From the Calculations}
The calculations are used here as a selection tool, not as the project end point. The lowest-LCOW literature route is \textbf{Absorption} (SEC = 0.38~kWh/kg, LCOW = 13.18~USD/m$^3$). For the buildable prototype package, the selected configuration is \textbf{a modular vapor-compression condensation skid with polishing treatment and buffer storage} because its LCOW penalty relative to the numerical minimum is only 0.80~USD/m$^3$, while it can be assembled from standard compressor, heat-exchanger, fan, tank, filter, UV, and pump hardware.
\paragraph{Calculated Best-Option Parameters}
\begin{itemize}[leftmargin=*]
\item Full-scale reference block: $P_{\mathrm{dc}}=200$~MW, $\eta_{\mathrm{mod}}=0.215$, $A_{\mathrm{pv}}=930,233$~m$^2$.
\item Cleaning basis: $w=1.8$~L/m$^2$/event, $N_{\mathrm{clean}}=18$/yr, $f_{\mathrm{margin}}=0.25$.
\item Required water output: $V_{\mathrm{year}}=37,674.4$~m$^3$/yr and $m_{\mathrm{day,avg}}=103,218$~kg/day ($\approx 103.2$~m$^3$/day).
\item Prototype-selected route: VCR SEC = 0.42~kWh/kg, $E_{\mathrm{year}}=15.82$~GWh/yr, LCOW = 13.98~USD/m$^3$.
\end{itemize}
\paragraph{Traceability Matrix}
\begin{tabular}{@{}p{3.6cm}p{4.2cm}p{7.0cm}@{}}
\toprule
Calculation output & Design decision & Build implication \\
\midrule
$A_{\mathrm{pv}}=930,233$~m$^2$ & Storage and production sizing & Sets the full-scale water duty and the prototype scale ratio. \\
$V_{\mathrm{year}}=37,674.4$~m$^3$/yr & Tank and pump duty & Fixes average delivery need of 103.2~m$^3$/day. \\
Lowest LCOW = 13.18~USD/m$^3$ & Screen candidate routes & Keeps absorption and VCR as front-runners. \\
VCR LCOW gap = 0.80~USD/m$^3$ & Select prototype route & Accepts small cost penalty to gain procurement and assembly readiness. \\
\bottomrule
\end{tabular}
\paragraph{Simulation Plan}
\begin{enumerate}[leftmargin=*]
\item Use hourly NASA POWER and ERA5 weather to build realistic temperature and humidity boundary conditions.
\item Validate the PV block in PVsyst with the same soiling schedule that justified the cleaning frequency.
\item Simulate the VCR prototype in Python with hourly inlet state, fan power, compressor cycling, and condensate accumulation.
\item Run sensitivity cases for summer peak temperature, low-RH nights, dust-loading on filters, and one failed-cleaning-cycle recovery.
\item Accept the selected parameters only if condensate production meets the target for 95\% of simulated days and polished-water TDS stays within the rinse specification.
\end{enumerate}
\paragraph{Engineering Drawings, BOM, and Component Specifications}
\begin{itemize}[leftmargin=*]
\item Drawing register: D-101 process flow diagram, D-102 skid general arrangement (2200\,$\times$\,1400\,$\times$\,2100~mm), D-103 condensate tank and drain routing, D-104 treatment manifold and sampling points, D-105 electrical one-line and control panel layout, D-106 instrumentation loop diagram, and D-107 cleaning-water test manifold to the autonomous cleaning interface.
\item Prototype BOM: hermetic compressor package sized for about 200~L/day condensate, hydrophilic evaporator coil, air-cooled condenser, two EC axial fans, 500~L raw-water tank, 500~L polished-water tank, 5~$\mu$m sediment filter, activated-carbon filter, UV reactor, delivery pump rated 2.0~m$^3$/h at 3.0~bar, flow meter, TDS sensor, temperature/RH sensors, PLC or datalogger, isolation valves, relief devices, and stainless tubing.
\item Component specification targets: 220~V single-phase laboratory supply for the prototype, food-safe or HDPE water-contact surfaces, sealed electrical enclosure, and service access for filter replacement.
\end{itemize}
\paragraph{Build Manual, Safety, and Quality Control}
\begin{enumerate}[leftmargin=*]
\item Assemble the refrigeration loop and pressure-test it before energizing any electrical equipment.
\item Install condensate collection, treatment cartridges, UV unit, and final storage in the same sequence shown on D-101.
\item Calibrate all sensors, then execute dry commissioning, wet commissioning, and 24-hour condensate trials.
\item Apply lockout-tagout, refrigerant handling controls, splash protection, electrical earthing checks, and biological cleaning of tanks before each trial.
\item Hold quality checkpoints at leak test, insulation resistance, airflow verification, condensate rate verification, TDS verification, and cleaning-interface pressure verification.
\end{enumerate}
\paragraph{Integrated Master Project Plan}
\begin{tabular}{@{}p{2.2cm}p{4.4cm}p{3.2cm}p{2.4cm}p{2.8cm}@{}}
\toprule
Milestone & Deliverable & Resources & Date & Success criteria \\
\midrule
M1 & Freeze design basis and traceability matrix & PI, analyst & T0+2d & YAML, formulas, and route ranking approved \\
M2 & Complete PVsyst and hourly AWH simulations & PV analyst, controls student & T0+7d & Weather-linked performance model closed \\
M3 & Release drawing package D-101 to D-107 & CAD lead, mechanical lead & T0+10d & All interfaces dimensioned and reviewed \\
M4 & Procure prototype BOM and fabricate skid & Procurement, workshop technician & T0+17d & All long-lead items received and inspected \\
M5 & Build and commission prototype & Mechanical, electrical, safety officer & T0+24d & 24-hour safe operation achieved \\
M6 & Validate against calculation targets & Test engineer, QA lead & T0+28d & Condensate rate, TDS, and pump duty meet target \\
\bottomrule
\end{tabular}
